\section{IP Addressing Schema (IPv4 and IPv6)}\label{sec:ip}
% Rubric criterion 3 (15 %): VLSM + IPv6 hierarchy, scalable to 10k+ endpoints, summarisation at the edges.

\subsection{Design rules}\label{sec:ip-rules}
The plan uses private IPv4 space (RFC~1918 \cite{rfc1918}) behind NAT and a global IPv6 hierarchy; the
IPv6 prefix \texttt{2001:db8::/32} is the documentation prefix \cite{rfc3849} and stands for the block TVI
would receive from its provider. Six rules govern every assignment:

\begin{enumerate}
  \item \textbf{One summarisable block per site.} HCMC \texttt{10.1.0.0/16}, Hanoi \texttt{10.2.0.0/16},
        Da Nang \texttt{10.3.0.0/16}; in IPv6 one /48 per site (\texttt{2001:db8:1::/48} \ldots).
        Each site announces exactly one prefix per address family to the rest of the company.
  \item \textbf{VLSM, largest block first, aligned to its size.} Subnets are sized for today's demand plus
        headroom. No user VLAN is larger than a /22 (1{,}022 hosts), which bounds the broadcast domain; only
        the passive, storm-controlled IoT VLANs in Da Nang are /20s.
  \item \textbf{Gateway convention.} In every subnet \texttt{.1} is the HSRP virtual gateway, \texttt{.2} and
        \texttt{.3} are the two physical gateways, \texttt{.4}--\texttt{.20} are reserved for static
        infrastructure and excluded from DHCP.
  \item \textbf{IPv6 mirrors IPv4.} One /64 per VLAN, with the VLAN ID as the fourth hextet: VLAN~20 in Hanoi
        is \texttt{2001:db8:2:20::/64}. An address can be read without a lookup table \cite{rfc4291}.
  \item \textbf{Infrastructure kept apart from users.} Loopbacks \texttt{10.<site>.255.0/24}, HQ routed links
        \texttt{10.1.250.0/24}, tunnels \texttt{10.255.1.0/24}, management VLAN \texttt{10.<site>.99.0/24}.
  \item \textbf{Provider space} on the WAN uses the documentation ranges of RFC~5737 \cite{rfc5737}:
        ISP-A \texttt{203.0.113.0/24}, ISP-B \texttt{198.51.100.0/24}.
\end{enumerate}

The complete plan, including the reserved blocks, is in \path{Design/IP_Schema.md}; the Excel version
\path{Design/IP_Schema.xlsx} recomputes the capacity figures with live formulas.

\subsection{Capacity and growth to 2030}\label{sec:ip-capacity}
\Cref{tab:sites} compares demand with addressable capacity. User capacity counts the usable addresses of
the user VLANs (executive, R\&D, AI lab, logistics); IoT capacity those of the IoT VLANs.

% GENERATED by Tools/gen_latex.py - do not edit; change Tools/tvi_model.py and regenerate.
\begin{table}[htbp]
\centering\scriptsize\setlength{\tabcolsep}{3pt}
\caption{Site summary: demand, addressable capacity and aggregates}\label{tab:sites}
\begin{tabularx}{\textwidth}{X r r r r r l l}
\toprule
\textbf{Site} & \textbf{Users now} & \textbf{Users 2030} & \textbf{User cap.} & \textbf{IoT cap.} & \textbf{Reserved} & \textbf{IPv4 aggregate} & \textbf{IPv6 aggregate} \\
\midrule
HCMC & 2000 & 4000 & 2551 & 1021 & 35840 & \texttt{10.1.0.0/16} & \texttt{2001:db8:1::/48} \\
HANOI & 1500 & 4000 & 2295 & 253 & 35840 & \texttt{10.2.0.0/16} & \texttt{2001:db8:2::/48} \\
DANANG & 500 & 2000 & 1335 & 12279 & 29696 & \texttt{10.3.0.0/16} & \texttt{2001:db8:3::/48} \\
\bottomrule
\end{tabularx}
\end{table}


Today's user VLANs cover today's head count with headroom at every site (HCMC 2{,}000 users on 2{,}551
addresses, Hanoi 1{,}500 on 2{,}295, Da Nang 500 on 1{,}335). The 2030 head count is covered by the reserve:
every /16 keeps between 29{,}696 and 35{,}840 addresses in contiguous blocks, enough for the next
generation of user VLANs without renumbering any existing subnet. The IoT requirement drives the Da Nang plan: three /20s for
sensors, warehouse automation and scanners/CCTV give \FactIotCapacityDanang{} usable addresses today, and
the reserved blocks \texttt{10.3.176.0/20}, \texttt{10.3.192.0/19} and \texttt{10.3.224.0/20} take the hub
beyond 30{,}000 IoT endpoints. HQ reserves \texttt{10.1.128.0/18} for the EVPN-VXLAN data-centre fabric
(\cref{sec:evpn}).

The sizing deliberately avoids two kinds of waste. Giving every VLAN a /20 ``to be safe'' would create
broadcast domains of 4{,}000 hosts and consume the reserve; giving them exactly today's size would force
renumbering within two years. A /22 for each large user group and contiguous reserves per site balance
the two.

\subsection{VLAN catalogue}\label{sec:ip-vlans}
VLAN IDs have the same meaning at every site, so a policy written for VLAN~50 is the guest policy
everywhere (\cref{tab:vlans}). VLAN~998 is the native VLAN of every trunk and carries no hosts; unused
access ports are shut and placed in VLAN~999.
% GENERATED by Tools/gen_latex.py - do not edit; change Tools/tvi_model.py and regenerate.
\begin{table}[htbp]
\centering\footnotesize
\caption{VLAN catalogue}\label{tab:vlans}
\begin{tabularx}{\textwidth}{r l X l}
\toprule
\textbf{VLAN} & \textbf{Name} & \textbf{Purpose} & \textbf{Sites} \\
\midrule
10 & \texttt{EXEC\_FINANCE} & Executives and Finance (maximum protection) & HCMC, HANOI, DANANG \\
20 & \texttt{RND\_ENGINEERING} & Engineering / general corporate staff & HCMC, HANOI, DANANG \\
21 & \texttt{AI\_LAB\_COMPUTE} & AI-lab compute and private-cloud access (Hanoi, east-west heavy) & HANOI \\
30 & \texttt{LOGISTICS\_OPS} & Logistics and warehouse operations staff & HCMC, DANANG \\
40 & \texttt{VOICE} & VoIP phones and video endpoints (EF/AF31/CS3) & HCMC, HANOI, DANANG \\
50 & \texttt{GUEST} & Guests and visitors (Internet only) & HCMC, HANOI, DANANG \\
60 & \texttt{IOT\_SENSORS} & IoT: trackers / RFID / environmental sensors & HCMC, HANOI, DANANG \\
61 & \texttt{IOT\_AUTOMATION} & IoT: AGVs, robots, PLC / warehouse automation & DANANG \\
62 & \texttt{IOT\_SCANNERS\_CCTV} & IoT: handheld scanners, CCTV & DANANG \\
90 & \texttt{L3\_TRANSIT} & Core-to-core L3 transit (OSPF adjacency between the two core switches) & HCMC \\
99 & \texttt{MANAGEMENT} & Network management (SSH, SNMP, syslog, automation host) & HCMC, HANOI, DANANG \\
100 & \texttt{DC\_SERVERS} & Data-centre services: DHCP, DNS, NTP & HCMC \\
998 & \texttt{NATIVE\_UNUSED} & Unused native VLAN for trunks (VLAN-hopping defence) & all switches \\
999 & \texttt{BLACKHOLE} & Blackhole VLAN for unused / shut-down access ports & all switches \\
\bottomrule
\end{tabularx}
\end{table}


\subsection{Per-site VLSM plans}\label{sec:ip-vlsm}
HCMC (\cref{tab:vlsm-hq}) carries the largest user groups and the data-centre VLAN; VLAN~90 is the /30 transit
between the two cores. Hanoi (\cref{tab:vlsm-hn}) adds VLAN~21 for the AI-lab compute nodes, separated
from the R\&D desktops so that heavy east-west traffic and its security policy can be handled on its own.
Da Nang (\cref{tab:vlsm-dn}) has small office VLANs and the three IoT /20s.

% GENERATED by Tools/gen_latex.py - do not edit; change Tools/tvi_model.py and regenerate.
\begin{table}[htbp]
\centering\scriptsize\setlength{\tabcolsep}{4pt}
\caption{Ho Chi Minh City: IPv4 VLSM plan and IPv6 /64 per VLAN}\label{tab:vlsm-hq}
\begin{tabularx}{\textwidth}{@{}r l l r l X l@{}}
\toprule
\textbf{VLAN} & \textbf{Name} & \textbf{IPv4 subnet} & \textbf{Usable} & \textbf{Gateway (VIP)} & \textbf{DHCP range} & \textbf{IPv6 /64} \\
\midrule
10 & \texttt{EXEC\_FINANCE} & \texttt{10.1.0.0/23} & 510 & \texttt{10.1.0.1} & \texttt{10.1.0.21 - 10.1.1.254} & \texttt{2001:db8:1:10::/64} \\
20 & \texttt{RND\_ENGINEERING} & \texttt{10.1.4.0/22} & 1\,022 & \texttt{10.1.4.1} & \texttt{10.1.4.21 - 10.1.7.254} & \texttt{2001:db8:1:20::/64} \\
30 & \texttt{LOGISTICS\_OPS} & \texttt{10.1.8.0/22} & 1\,022 & \texttt{10.1.8.1} & \texttt{10.1.8.21 - 10.1.11.254} & \texttt{2001:db8:1:30::/64} \\
40 & \texttt{VOICE} & \texttt{10.1.12.0/22} & 1\,022 & \texttt{10.1.12.1} & \texttt{10.1.12.21 - 10.1.15.254} & \texttt{2001:db8:1:40::/64} \\
50 & \texttt{GUEST} & \texttt{10.1.2.0/23} & 510 & \texttt{10.1.2.1} & \texttt{10.1.2.21 - 10.1.3.254} & \texttt{2001:db8:1:50::/64} \\
60 & \texttt{IOT\_SENSORS} & \texttt{10.1.16.0/22} & 1\,022 & \texttt{10.1.16.1} & \texttt{10.1.16.21 - 10.1.19.254} & \texttt{2001:db8:1:60::/64} \\
90 & \texttt{L3\_TRANSIT} & \texttt{10.1.248.0/30} & 2 & \texttt{-} & - & \texttt{2001:db8:1:90::/64} \\
99 & \texttt{MANAGEMENT} & \texttt{10.1.99.0/24} & 254 & \texttt{10.1.99.1} & static & \texttt{2001:db8:1:99::/64} \\
100 & \texttt{DC\_SERVERS} & \texttt{10.1.100.0/24} & 254 & \texttt{10.1.100.1} & static & \texttt{2001:db8:1:100::/64} \\
\bottomrule
\end{tabularx}
\end{table}

% GENERATED by Tools/gen_latex.py - do not edit; change Tools/tvi_model.py and regenerate.
\begin{table}[htbp]
\centering\scriptsize\setlength{\tabcolsep}{4pt}
\caption{Hanoi (R\&D campus): IPv4 VLSM plan and IPv6 /64 per VLAN}\label{tab:vlsm-hn}
\begin{tabularx}{\textwidth}{@{}r l l r l X l@{}}
\toprule
\textbf{VLAN} & \textbf{Name} & \textbf{IPv4 subnet} & \textbf{Usable} & \textbf{Gateway (VIP)} & \textbf{DHCP range} & \textbf{IPv6 /64} \\
\midrule
10 & \texttt{EXEC\_FINANCE} & \texttt{10.2.0.0/24} & 254 & \texttt{10.2.0.1} & \texttt{10.2.0.21 - 10.2.0.254} & \texttt{2001:db8:2:10::/64} \\
20 & \texttt{RND\_ENGINEERING} & \texttt{10.2.4.0/22} & 1\,022 & \texttt{10.2.4.1} & \texttt{10.2.4.21 - 10.2.7.254} & \texttt{2001:db8:2:20::/64} \\
21 & \texttt{AI\_LAB\_COMPUTE} & \texttt{10.2.8.0/22} & 1\,022 & \texttt{10.2.8.1} & \texttt{10.2.8.21 - 10.2.11.254} & \texttt{2001:db8:2:21::/64} \\
40 & \texttt{VOICE} & \texttt{10.2.12.0/22} & 1\,022 & \texttt{10.2.12.1} & \texttt{10.2.12.21 - 10.2.15.254} & \texttt{2001:db8:2:40::/64} \\
50 & \texttt{GUEST} & \texttt{10.2.1.0/24} & 254 & \texttt{10.2.1.1} & \texttt{10.2.1.21 - 10.2.1.254} & \texttt{2001:db8:2:50::/64} \\
60 & \texttt{IOT\_SENSORS} & \texttt{10.2.16.0/24} & 254 & \texttt{10.2.16.1} & \texttt{10.2.16.21 - 10.2.16.254} & \texttt{2001:db8:2:60::/64} \\
99 & \texttt{MANAGEMENT} & \texttt{10.2.99.0/24} & 254 & \texttt{10.2.99.1} & static & \texttt{2001:db8:2:99::/64} \\
\bottomrule
\end{tabularx}
\end{table}

% GENERATED by Tools/gen_latex.py - do not edit; change Tools/tvi_model.py and regenerate.
\begin{table}[htbp]
\centering\scriptsize\setlength{\tabcolsep}{4pt}
\caption{Da Nang (logistics hub): IPv4 VLSM plan and IPv6 /64 per VLAN}\label{tab:vlsm-dn}
\begin{tabularx}{\textwidth}{@{}r l l r l X l@{}}
\toprule
\textbf{VLAN} & \textbf{Name} & \textbf{IPv4 subnet} & \textbf{Usable} & \textbf{Gateway (VIP)} & \textbf{DHCP range} & \textbf{IPv6 /64} \\
\midrule
10 & \texttt{EXEC\_FINANCE} & \texttt{10.3.0.0/26} & 62 & \texttt{10.3.0.1} & \texttt{10.3.0.21 - 10.3.0.62} & \texttt{2001:db8:3:10::/64} \\
20 & \texttt{RND\_ENGINEERING} & \texttt{10.3.1.0/24} & 254 & \texttt{10.3.1.1} & \texttt{10.3.1.21 - 10.3.1.254} & \texttt{2001:db8:3:20::/64} \\
30 & \texttt{LOGISTICS\_OPS} & \texttt{10.3.4.0/22} & 1\,022 & \texttt{10.3.4.1} & \texttt{10.3.4.21 - 10.3.7.254} & \texttt{2001:db8:3:30::/64} \\
40 & \texttt{VOICE} & \texttt{10.3.8.0/23} & 510 & \texttt{10.3.8.1} & \texttt{10.3.8.21 - 10.3.9.254} & \texttt{2001:db8:3:40::/64} \\
50 & \texttt{GUEST} & \texttt{10.3.0.128/25} & 126 & \texttt{10.3.0.129} & \texttt{10.3.0.149 - 10.3.0.254} & \texttt{2001:db8:3:50::/64} \\
60 & \texttt{IOT\_SENSORS} & \texttt{10.3.128.0/20} & 4\,094 & \texttt{10.3.128.1} & \texttt{10.3.128.21 - 10.3.143.254} & \texttt{2001:db8:3:60::/64} \\
61 & \texttt{IOT\_AUTOMATION} & \texttt{10.3.144.0/20} & 4\,094 & \texttt{10.3.144.1} & \texttt{10.3.144.21 - 10.3.159.254} & \texttt{2001:db8:3:61::/64} \\
62 & \texttt{IOT\_SCANNERS\_CCTV} & \texttt{10.3.160.0/20} & 4\,094 & \texttt{10.3.160.1} & \texttt{10.3.160.21 - 10.3.175.254} & \texttt{2001:db8:3:62::/64} \\
99 & \texttt{MANAGEMENT} & \texttt{10.3.99.0/24} & 254 & \texttt{10.3.99.1} & static & \texttt{2001:db8:3:99::/64} \\
\bottomrule
\end{tabularx}
\end{table}


\subsection{Infrastructure addressing and summarisation}\label{sec:ip-infra}
\Cref{tab:links} lists every point-to-point link, tunnel and provider hand-off. Each site's edge routers
install a static route for the site aggregate to \texttt{Null0} (for example
\texttt{ip route 10.2.0.0 255.255.0.0 Null0}). This route has two jobs: BGP originates a \texttt{network}
statement only if exactly that prefix is in the routing table, and packets for unused parts of the block
are discarded at the edge instead of being sent back towards the provider. As a result, the inter-site BGP
table holds three IPv4 prefixes and three IPv6 prefixes however many VLANs are added -- a VLAN change in one
site never reaches the routing tables of the others.

\begin{landscape}
% GENERATED by Tools/gen_latex.py - do not edit; change Tools/tvi_model.py and regenerate.
\scriptsize
\begin{longtable}{l l l l l l}
\caption{Point-to-point links, overlay tunnels and provider hand-offs}\label{tab:links}\\
\toprule
\textbf{Type} & \textbf{Device A / interface} & \textbf{IPv4 A} & \textbf{Device B / interface} & \textbf{IPv4 B} & \textbf{Subnet} \\
\midrule
\endfirsthead
\toprule
\textbf{Type} & \textbf{Device A / interface} & \textbf{IPv4 A} & \textbf{Device B / interface} & \textbf{IPv4 B} & \textbf{Subnet} \\
\midrule
\endhead
\bottomrule
\endfoot
Core-Edge p2p & HCMC-CORE-SW01 Gi1/0 & \texttt{10.1.250.1} & HCMC-EDGE-RTR01 Gi0/1 & \texttt{10.1.250.2} & \texttt{10.1.250.0/30} \\
Core-Edge p2p & HCMC-CORE-SW01 Gi1/1 & \texttt{10.1.250.5} & HCMC-EDGE-RTR02 Gi0/1 & \texttt{10.1.250.6} & \texttt{10.1.250.4/30} \\
Core-Edge p2p & HCMC-CORE-SW02 Gi1/0 & \texttt{10.1.250.9} & HCMC-EDGE-RTR01 Gi0/2 & \texttt{10.1.250.10} & \texttt{10.1.250.8/30} \\
Core-Edge p2p & HCMC-CORE-SW02 Gi1/1 & \texttt{10.1.250.13} & HCMC-EDGE-RTR02 Gi0/2 & \texttt{10.1.250.14} & \texttt{10.1.250.12/30} \\
Edge-Edge p2p & HCMC-EDGE-RTR01 Gi0/3 & \texttt{10.1.250.17} & HCMC-EDGE-RTR02 Gi0/3 & \texttt{10.1.250.18} & \texttt{10.1.250.16/30} \\
GRE/IPsec Tunnel1 (ISP-A) & HCMC-EDGE-RTR01 Tunnel1 & \texttt{10.255.1.1} & HANOI-EDGE-RTR01 Tunnel1 & \texttt{10.255.1.2} & \texttt{10.255.1.0/30} \\
GRE/IPsec Tunnel2 (ISP-B) & HCMC-EDGE-RTR02 Tunnel2 & \texttt{10.255.1.5} & HANOI-EDGE-RTR02 Tunnel2 & \texttt{10.255.1.6} & \texttt{10.255.1.4/30} \\
GRE/IPsec Tunnel3 (ISP-A) & HCMC-EDGE-RTR01 Tunnel3 & \texttt{10.255.1.9} & DANANG-EDGE-RTR01 Tunnel3 & \texttt{10.255.1.10} & \texttt{10.255.1.8/30} \\
GRE/IPsec Tunnel4 (ISP-B) & HCMC-EDGE-RTR02 Tunnel4 & \texttt{10.255.1.13} & DANANG-EDGE-RTR02 Tunnel4 & \texttt{10.255.1.14} & \texttt{10.255.1.12/30} \\
WAN to ISP-A-RTR & HCMC-EDGE-RTR01 Gi0/0 & \texttt{203.0.113.2} & ISP-A-RTR Gi0/0 & \texttt{203.0.113.1} & \texttt{203.0.113.0/30} \\
WAN to ISP-A-RTR & HANOI-EDGE-RTR01 Gi0/0 & \texttt{203.0.113.6} & ISP-A-RTR Gi0/1 & \texttt{203.0.113.5} & \texttt{203.0.113.4/30} \\
WAN to ISP-A-RTR & DANANG-EDGE-RTR01 Gi0/0 & \texttt{203.0.113.10} & ISP-A-RTR Gi0/2 & \texttt{203.0.113.9} & \texttt{203.0.113.8/30} \\
WAN to ISP-B-RTR & HCMC-EDGE-RTR02 Gi0/0 & \texttt{198.51.100.2} & ISP-B-RTR Gi0/0 & \texttt{198.51.100.1} & \texttt{198.51.100.0/30} \\
WAN to ISP-B-RTR & HANOI-EDGE-RTR02 Gi0/0 & \texttt{198.51.100.6} & ISP-B-RTR Gi0/1 & \texttt{198.51.100.5} & \texttt{198.51.100.4/30} \\
WAN to ISP-B-RTR & DANANG-EDGE-RTR02 Gi0/0 & \texttt{198.51.100.10} & ISP-B-RTR Gi0/2 & \texttt{198.51.100.9} & \texttt{198.51.100.8/30} \\
ISP transit & ISP-A-RTR Gi0/3 & \texttt{203.0.113.13} & INTERNET-RTR Gi0/0 & \texttt{203.0.113.14} & \texttt{203.0.113.12/30} \\
ISP transit & ISP-B-RTR Gi0/3 & \texttt{198.51.100.13} & INTERNET-RTR Gi0/1 & \texttt{198.51.100.14} & \texttt{198.51.100.12/30} \\
\end{longtable}
\normalsize

\end{landscape}
